% !PRINTED 279
\weberpage{279}
%<<CONT
zu setzen, woraus sich durch Bildung von $R\,A\,R^{-1}$ die erzeugenden Substitutionen der Octaëdergruppe in reeller Form ergeben:
\[
\Theta = \begin{pmatrix} 2, & 1 \\ 3, & 2 \end{pmatrix}, \
\psi = \begin{pmatrix} 2, & 2 \\ 1, & -2 \end{pmatrix}, \
\omega = \psi\,\Theta = \begin{pmatrix} 3, & -1 \\ 3, & -3 \end{pmatrix}
\]
\weq{11}{
\chi = \begin{pmatrix} 0, & 2 \\ 3, & 1 \end{pmatrix}, \
\chi^2 = \begin{pmatrix} -1, & 2 \\ 3, & 0 \end{pmatrix}.
}
Nun ist es leicht, nach §. 57 die vollständige Gruppe zu bilden, was nicht nöthig ist, weiter auszuführen.

\gesperrt{Die Octaëdergruppe für $p=7$ ist ein Theiler von $L_7$ vom Index 7.}

V. \gesperrt{Ikosaëdergruppe.}

In dem Falle $p=5$ ist $L_p$ selbst eine Ikosaëdergruppe. Für andere Werthe von $p$ kann nur dann eine Ikosaëdergruppe in $L_p$ enthalten sein, wenn $p^2-1$ durch 5 theilbar, also $p\equiv\pm 1\ (\mathrm{mod}\ 5)$ ist. Dann erhalten wir aus §. 58 die erzeugenden Substitutionen einer solchen Gruppe.

Setzen wir zur Abkürzung
\[
\gamma^{\frac{p^2-1}{10}} = \varrho,
\]
so haben wir $\varrho^2$ für $\varepsilon$ in die Formeln des §. 58 einzusetzen, und erhalten
\weq{12}{
\Theta = \begin{pmatrix} \varrho, & 0 \\ 0, & \varrho^{-1} \end{pmatrix}, \
\psi = \begin{pmatrix} 0, & 1 \\ -1, & 0 \end{pmatrix}, \quad
\chi = \begin{pmatrix} \dfrac{-1}{\varrho^2-\varrho^{-2}}, & \dfrac{1}{\varrho-\varrho^{-1}} \\[2mm] \dfrac{1}{\varrho-\varrho^{-1}}, & \dfrac{1}{\varrho^2-\varrho^{-2}} \end{pmatrix}\ [\S.\,58,\,(26)].
}
Ist $p\equiv 1\ (\mathrm{mod}\ 5)$, so ist $\varrho$, und damit die ganze Gruppe reell. Ist aber $p\equiv -1\ (\mathrm{mod}\ 5)$, so sind $\varrho$ und $\varrho^{-1}$ conjugirt, und folglich ist die gefundene Gruppe in $\Gamma_p$ enthalten, und um die Ikosaëdergruppe in $L_p$ zu erhalten, müssen wir erst nach §. 69 transformiren.

\gesperrt{Für $p=11$ erhalten wir unmittelbar eine reelle Ikosaëdergruppe vom Index 11.}
